\chapter{Exercise Physiologist}

\section*{Definition and Overview}
An exercise physiologist is an allied health professional who uses exercise as a treatment for medical conditions and injuries. They design and supervise individualised exercise programs for people with heart disease, diabetes, lung disease, obesity, musculoskeletal conditions, mental health problems, and many other conditions. Exercise physiologists assess fitness and health, set goals, and monitor progress, helping people exercise safely and effectively. They work with GPs and other health professionals as part of the care team.

\section*{Epidemiology}
Exercise physiologists are an increasingly important part of Australian health care, reflecting the strong evidence that exercise treats and prevents chronic disease. They work with people with heart disease, diabetes, and lung disease, with people recovering from injury or surgery, and with those managing weight and mental health. Referrals have grown as Medicare and other funding arrangements recognise the value of exercise as medicine.

\section*{Aetiology and Pathophysiology}
Exercise physiologists apply the science of exercise to clinical care.
\begin{itemize}
    \item \textbf{Training:} they complete university study in exercise science and clinical exercise physiology, with supervised clinical practice
    \item \textbf{Assessment:} they evaluate fitness, strength, movement, and health before designing a program
    \item \textbf{The programs:} individualised exercise plans that address the specific condition and goals
    \item \textbf{How exercise works:} it strengthens the heart, improves blood sugar and blood pressure, builds muscle, improves movement, and supports mental health
    \item \textbf{Safety:} they understand the risks of exercise in illness and design programs that are safe
    \item \textbf{Accreditation:} accredited exercise physiologists are members of Exercise and Sports Science Australia
\end{itemize}

\section*{Clinical Features}
People are referred to an exercise physiologist for many reasons.
\begin{itemize}
    \item Heart disease, including after a heart attack or cardiac surgery
    \item Diabetes and metabolic disease
    \item Lung disease, including COPD
    \item Obesity and weight management
    \item Musculoskeletal conditions and injuries
    \item Neurological conditions, including stroke and Parkinson's disease
    \item Mental health conditions
    \item Cancer rehabilitation and post-treatment care
\end{itemize}

\section*{Differential Diagnosis}
Other professionals provide related services.
\begin{itemize}
    \item Physiotherapists focus on restoring movement and treating injuries
    \item Personal trainers provide general fitness training for healthy people
    \item Dietitians address nutrition alongside the exercise
    \item Occupational therapists focus on daily function
    \item The exercise physiologist specialises in exercise as treatment for medical conditions
\end{itemize}

\section*{When to Seek Medical Care}
People with chronic conditions, recovering from illness or surgery, or seeking help to become active should ask their GP about referral to an exercise physiologist. Medicare rebates are available for some people with chronic conditions.

\textbf{Call 000 (or local emergency services) for:}
\begin{itemize}
    \item Chest pain, collapse, or severe breathlessness during exercise
    \item Severe dizziness or palpitations
    \item Sudden weakness or difficulty speaking
    \item Any exercise-related emergency
\end{itemize}

\section*{Diagnosis and Investigations}
The exercise physiologist begins with a thorough assessment.
\begin{itemize}
    \item \textbf{Health history:} conditions, medicines, and goals
    \item \textbf{Fitness assessment:} aerobic fitness, strength, and mobility testing
    \item \textbf{Risk assessment:} understanding the risks of exercise for the individual
    \item \textbf{Goal setting:} specific, achievable goals
    \item \textbf{Program design:} an individualised plan
    \item \textbf{Monitoring:} regular review and adjustment
\end{itemize}

\section*{Management}
The exercise physiologist provides structured exercise care.
\begin{itemize}
    \item \textbf{Supervised sessions:} individual or group exercise sessions
    \item \textbf{Home programs:} exercise to do between sessions
    \item \textbf{Education:} understanding how exercise helps the condition
    \item \textbf{Progress tracking:} fitness and health measures
    \item \textbf{Coordination with GPs:} regular updates and review
    \item \textbf{Long-term support:} building sustainable exercise habits
\end{itemize}

\section*{Complications}
\begin{itemize}
    \item Injury if programs are not individualised
    \item Worsening of symptoms if exercise is inappropriate
    \item The risks of inactivity, which are far greater
    \item Cost and access barriers for some people
\end{itemize}

\section*{Prognosis and Outcomes}
Exercise programs delivered by exercise physiologists are effective across many conditions, improving fitness, strength, blood sugar, blood pressure, mood, and quality of life. People who complete structured programs gain the skills and confidence to exercise independently, and the benefits persist with ongoing activity.

\section*{Prevention and Self-Care}
\begin{itemize}
    \item Ask your GP about a referral, including the Medicare chronic disease management plan
    \item Attend sessions regularly and complete the program
    \item Follow the home exercise plan
    \item Ask questions and share any symptoms
    \item Track progress and celebrate achievements
    \item Continue activity after the program ends
\end{itemize}

\section*{Sources and Further Reading}
The following reputable sources provide further information for healthcare students and patients seeking reliable, current advice about exercise physiologists.
\begin{itemize}
    \item Exercise and Sports Science Australia. \textit{Accredited exercise physiologists.} https://www.essa.org.au/
    \item Healthdirect Australia. \textit{Exercise physiologist.} https://www.healthdirect.gov.au/exercise-physiologist
    \item Services Australia. \textit{Chronic disease management plans.} https://www.servicesaustralia.gov.au/
    \item Heart Foundation Australia. \textit{Cardiac rehabilitation.} https://www.heartfoundation.org.au/
    \item Diabetes Australia. \textit{Exercise and diabetes.} https://www.diabetesaustralia.com.au/
    \item Better Health Channel. \textit{Allied health services.} https://www.betterhealth.vic.gov.au/
\end{itemize}
